\documentclass[11pt,a4paper]{article}
\usepackage[T1]{fontenc}
\usepackage{lmodern,amsmath,amssymb,amsthm,mathtools}
\usepackage[margin=25mm]{geometry}
\usepackage{microtype,booktabs,enumitem,xcolor,listings,xurl,hyperref,fancyhdr}
\hypersetup{colorlinks=true,linkcolor=blue!45!black,citecolor=blue!45!black,urlcolor=blue!45!black,pdftitle={Prime-Power Divisibility in Sun's A179524 Conjecture}}
\setlength{\headheight}{14pt}
\pagestyle{fancy}
\fancyhf{}
\fancyhead[L]{Sun's A179524 divisibility conjecture}
\fancyhead[R]{Prime-power results}
\fancyfoot[C]{\thepage}
\newtheorem{theorem}{Theorem}[section]
\newtheorem{lemma}[theorem]{Lemma}
\newtheorem{corollary}[theorem]{Corollary}
\newtheorem{proposition}[theorem]{Proposition}
\theoremstyle{remark}\newtheorem{remark}[theorem]{Remark}
\newcommand{\Z}{\mathbb Z}
\newcommand{\F}{\mathbb F}
\newcommand{\Q}{\mathbb Q}
\newcommand{\CT}{\operatorname{CT}}
\newcommand{\rad}{\operatorname{rad}}
\newcommand{\vp}{v_p}
\lstset{language=Python,basicstyle=\ttfamily\footnotesize,keywordstyle=\color{blue!55!black},commentstyle=\color{gray!80!black},stringstyle=\color{gray!80!black},breaklines=true,columns=fullflexible,keepspaces=true,showstringspaces=false,frame=single,rulecolor=\color{gray!40},aboveskip=10pt,belowskip=10pt}

\title{\textbf{Prime-Power Divisibility in Sun's\\A179524 Conjecture}\\[5pt]
\large An integral analytic digit operator and a Catalan diagonal reduction}
\author{Research note prepared with ChatGPT for Keio}
\date{8 November 2026}
\begin{document}
\maketitle
\begin{abstract}
For \(a_n=\sum_k(-4)^k\binom nk^2\binom{n-k}k^2\), put
\(S(n)=\sum_{k<n}(20k+17)a_k\).
We prove \(p^e\mid S(p^em)\) for all \(e,m\ge1\) whenever
\(p=5\) or \(p\equiv11,13,17,19\pmod{20}\), and retain the exact
identity \(v_2(S(n))=v_2(n)\).
An exact digit operator on restricted integral \(p\)-adic power series,
with all carry boundaries included, permits simultaneous induction
for arbitrary analytic weights and three distinguished moments.
A separate one-moment cancellation treats \(p=5\).
For every odd prime, the same operator propagates any finite chain of
verified prime-power bases to all cofactors; consequently the full
conjecture is equivalent to its prime-power instances.
We also prove that the off-diagonal condition in the supplied
negative-index manuscript follows from its diagonal condition, by an
integral Catalan coordinate change.
An exact telescoping certificate proves the previously guessed
order-four recurrence. A fresh exact run through \(50000\), combined
with propagation, gives an explicit proved domain of natural density
greater than \(0.9992\).
The remaining split-prime-power bases are not proved.
No historical-priority claim is made for the general analytic lifting
strategy or for the target-specific conclusions.
\end{abstract}

\section{The unchanged target and the enlarged proved domain}

Use zero for a binomial coefficient with an inadmissible lower index.
Define
\begin{equation}\label{def}
 \begin{aligned}
 M(n,k)&=\binom nk\binom{n-k}k,\qquad
 a_n=\sum_{0\le2k\le n}(-4)^kM(n,k)^2,\\
 S(n)&=\sum_{k=0}^{n-1}(20k+17)a_k.
 \end{aligned}
\end{equation}
For \(0\le2k\le n\), \(M(n,k)=n!/[k!^2(n-2k)!]\).
The initial values are \(1,1,-15,-143\).
The original question is
\begin{equation}\label{target}
                         n\mid S(n)\qquad(n\ge1).
\end{equation}
It is Sun's Conjecture 5.8(ii) in the 2010 source
\cite[printed p.~33]{Sun2010}, with \(a_n=W_n(1,-4)\).
OEIS A179524 attributes it to 1 July 2010 \cite{OEIS}.
The numbered source is version 11; the final arXiv version was
withdrawn to merge material into other papers.

The supplied manuscript \cite{Supplied} contributes an exact
negative-index continuation and two-index formulations. The present
note combines them with the previous CM prime-block argument, develops
an exact analytic operator, proves propagation through all layers,
and supplies a symbolic certificate for its order-four recurrence.
All arguments needed for the stated results are included here.

Let
\[
 \mathcal G=\{5\}\cup
 \{p\text{ prime}:p\equiv11,13,17,19\pmod{20}\},\qquad N_0=50000.
\]
For each odd prime \(p\), put
\[
 e_p=
 \begin{cases}
 \infty,&p\in\mathcal G,\\
 \max\bigl(\{1\}\cup\{e\ge1:p^e\le N_0\}\bigr),&p\notin\mathcal G.
 \end{cases}
\]
The minimum of a finite integer and \(\infty\) is that integer.

\begin{theorem}[Combined divisibility result]\label{main}
For every \(n\ge1\),
\[
 \boxed{\quad
 v_2(S(n))=v_2(n),\qquad
 2^{v_2(n)}
 \prod_{\substack{p\mid n\\p\ {\rm odd}}}
 p^{\min(v_p(n),e_p)}\ \mid S(n).
 \quad}
\]
In particular, \eqref{target} holds for every \(n\) with
\(v_p(n)\le e_p\) at each remaining odd prime. There is no exponent
restriction at \(2\) or at any prime in \(\mathcal G\).
\end{theorem}

The unrestricted prime families are proved analytically in
Corollaries~\ref{an:inert} and~\ref{five:all}.
The finite bottom chains are certified in Appendix~\ref{app:verify}
using the recurrence proved in Appendix~\ref{rec:section}, and then
propagated to every cofactor by Theorem~\ref{layer:prop}.
The exact \(2\)-part and prime cases are proved below.

The remaining finite exponent bounds are particularly concrete:
\[
\begin{array}{c|c}
\text{remaining prime}&\text{all-cofactor exponent proved}\\ \hline
3&9\\
7&5\\
23,29&3\\
41,43,47,61,67,83,89,101,103,107,109&2\\
127,149,163,167,181,223&2\\
p>223,\ p\equiv1,3,7,9\pmod{20}&1
\end{array}
\]
For example, \(3^9\mid S(3^9m)\) and \(7^5\mid S(7^5m)\)
hold for every positive \(m\). These are not finite-cofactor tests.

\begin{corollary}[A rigorous density bound]\label{density}
There is an explicit set of integers of natural density
greater than \(0.9992\) on which \eqref{target} holds.
\end{corollary}
\begin{proof}
Take the smaller set with arbitrary \(v_2(n)\); at odd primes
\(p\le223\), allow arbitrary exponent if \(p\in\mathcal G\), and exponent
at most \(e_p\) otherwise; at every prime \(p>223\), allow exponent
at most one. Define
\[
 \nu_p=
 \begin{cases}
 1,&p\in\mathcal G,\\
 1-p^{-(e_p+1)},&p\notin\mathcal G.
 \end{cases}
\]
The natural density of this set is
\[
 \delta_0=\frac8{\pi^2}
       \prod_{\substack{3\le p\le223\\p\ {\rm prime}}}
       \frac{\nu_p}{1-p^{-2}}
       =0.99921205018052892427\ldots .
\]
Existence follows by the Chinese remainder theorem for finitely many
local conditions and the tail estimate \(\sum_{p>y}p^{-2}\to0\).
The rigorous inequality \(\delta_0>0.9992\) follows using
\(\pi<355/113\) and exact rational arithmetic for the finite product;
that rational comparison is in Appendix~\ref{app:verify}.
The larger domain in Theorem~\ref{main} also allows unrestricted
exponents at inert primes above \(223\).
This density measures a proved integer domain, not a proportion of
the proof difficulty resolved.
\end{proof}
\section{Prime divisibility and Lucas propagation}

We first record two integer identities:
\begin{align}
 \binom jm^2
 &=\sum_{i=0}^m\binom mi\binom{m+i}i\binom j{m+i},
 \label{product}\\
 \sum_{j<n}(20j+17)\binom jt
 &=(20t+17)\binom n{t+1}
       +20(t+1)\binom n{t+2}. \label{hockey}
\end{align}
The first counts ordered pairs of \(m\)-subsets by their union size.
The second follows from the hockey-stick identity and
\(j\binom jt=t\binom jt+(t+1)\binom j{t+1}\).

\begin{lemma}[Ramanujan reduction]\label{primelemma}
For an odd prime \(p\), \(h=(p-1)/2\), and \(d=\lfloor p/4\rfloor\),
\begin{equation}\label{ramred}
 S(p)\equiv-(-1)^h\sum_{j=0}^{d}
       (20j+3)\frac{(4j)!}{j!^4(-1024)^j}\pmod p .
\end{equation}
\end{lemma}
\begin{proof}
For \(0\le k\le h\), the polynomial \(F_k(X)=\binom X{2k}^2\)
is defined over \(\F_p\) and satisfies
\(F_k(2k-1-X)=F_k(X)\). Thus
\[
 \sum_{r\in\F_p}(20r+17)F_k(r)
 =(20k+7)\sum_{r\in\F_p}F_k(r).
\]
Apply \eqref{product} and sum over \(0\le r<p\). Modulo \(p\), the
only surviving binomial \(\binom p{2k+i+1}\) has \(2k+i=p-1\).
Since \(i=p-1-2k\) is even, its coefficient reduces to
\(\binom{2k}{p-1-2k}\). Therefore
\[
 S(p)\equiv\sum_{k=\lceil h/2\rceil}^{h}
 (20k+7)(-4)^k\binom{2k}k^2\binom{2k}{2h-2k}\pmod p .
\]
Put \(k=h-j\). The needed congruences are
\[
 20k+7\equiv-(20j+3),\quad
 (-4)^{h-j}\equiv(-1)^h(-4)^{-j},
\]
\[
 \binom{p-1-2j}{h-j}^{\!2}
 \equiv\binom{h+j}{2j}^{\!2}
 \equiv\frac{\binom{2j}j^2}{256^j},\qquad
 \binom{p-1-2j}{2j}\equiv\binom{4j}{2j}.
\]
For the middle identity use
\(\binom{j-1/2}{2j}=(-1)^j\binom{2j}j/16^j\).
All denominators are units modulo \(p\). These formulas give
\eqref{ramred}.
\end{proof}

Sun's Theorem 1.2 \cite{SunRefinement} states, for \(p>3\),
\[
 \sum_{j=0}^{h}(20j+3)\frac{(4j)!}{j!^4(-1024)^j}
 \equiv(-1)^h p\bigl(2^{p-1}+2-(2^{p-1}-1)^2\bigr)\pmod {p^4}.
\]
For \(d<j\le h\), the numerator \((4j)!\) is divisible by \(p\)
and \(j!\) is not. Hence Lemma~\ref{primelemma} gives \(p\mid S(p)\).
The remaining primes satisfy \(S(2)=54\) and \(S(3)=-801\).

\begin{lemma}[Digit congruence]\label{lucas}
For every odd prime \(p\), \(u\ge0\), and \(0\le r<p\),
\[
                 a_{pu+r}\equiv a_u a_r\pmod p.
\]
\end{lemma}
\begin{proof}
Write the inner index as \(k=pv+s\), \(0\le s<p\).
Multinomial Lucas makes \(M(pu+r,pv+s)\) vanish modulo \(p\)
unless \(2s\le r\). In the latter case
\[
 M(pu+r,pv+s)\equiv M(u,v)M(r,s)\pmod p.
\]
The high-digit range is exactly \(0\le2v\le u\).
Also \((-4)^{pv+s}\equiv(-4)^{v+s}\pmod p\).
Summing the factored terms proves the claim.
\end{proof}

For every \(m\ge1\), it follows that
\begin{equation}\label{primeprop}
 S(pm)\equiv S(p)\sum_{u<m}a_u\equiv0\pmod p.
\end{equation}
The case \(p=2\) follows directly from \(a_u\equiv1\pmod2\).
Consequently \(\rad(n)\mid S(n)\) for every \(n\).
This argument alone does not increase the prime-power exponent.

\section{The entire 2-part, with exact valuation}

\begin{theorem}\label{twoadic}
For every \(n\ge1\), \(S(n)/n\) is an odd \(2\)-adic integer.
Thus \(S(n)\ne0\) and \(v_2(S(n))=v_2(n)\).
\end{theorem}
\begin{proof}
Rewrite \(M(j,k)=\binom j{2k}\binom{2k}k\).
For a fixed \(k\ge1\), use \eqref{product} with \(m=2k\), followed
by \eqref{hockey}, and divide by \(n\). Its contribution is an integer
linear combination of
\[
 (-4)^k\binom{2k}k^2
 \left(
 \frac{20t+17}{t+1}\binom{n-1}t+
 \frac{20(t+1)}{t+2}\binom{n-1}{t+1}
 \right),\qquad 2k\le t\le4k.
\]
The prefactor has \(2\)-adic valuation at least \(2k+2\), because
\(\binom{2k}k=2\binom{2k-1}{k-1}\).
Both displayed denominators are at most \(4k+2<2^{2k+2}\).
Every such term therefore has valuation at least one.
All sums are finite. The \(k=0\) contribution is
\[
                  \frac1n\sum_{j<n}(20j+17)=10n+7,
\]
which is odd. This proves the assertion, including \(n=1\).
\end{proof}

\section{A prime-square transfer valid for every cofactor}\label{sec:transfer}

Put \(f(n,k)=(-4)^kM(n,k)^2\), and define
\[
 L_n=\sum_k kf(n,k),\qquad
 U(N)=\sum_{n<N}a_n,\quad T(N)=\sum_{n<N}n a_n,\quad
 V(N)=\sum_{n<N}L_n.
\]
Thus \(S(N)=20T(N)+17U(N)\).
For an odd prime \(p\), let
\[
 \begin{aligned}
 Q_p&=\frac{(-4)^{p-1}-1}{p},\\
 A_r&=\sum_s f(r,s)(H_r-H_{r-2s}),\qquad
 B_r=\sum_s f(r,s)(H_{r-2s}-H_s)\quad(0\le r<p),
 \end{aligned}
\]
where \(H_0=0\), \(H_j=\sum_{i=1}^j1/i\), and only admissible \(s\)
are included. These harmonic numbers are \(p\)-integral.

\begin{lemma}[First-order digit expansion]\label{firstorder}
For every \(q\ge0\) and \(0\le r<p\),
\begin{equation}\label{digit2}
 a_{pq+r}\equiv a_q a_r+
 p\{2q a_q A_r+L_q(Q_p a_r+4B_r)\}\pmod {p^2}.
\end{equation}
\end{lemma}
\begin{proof}
Let \(E_p(N)=\prod_{1\le i\le N,\ p\nmid i}i\).
The exact factorization
\(N!=p^{\lfloor N/p\rfloor}\lfloor N/p\rfloor!E_p(N)\), together with
\(\sum_{i=1}^{p-1}i^{-1}=0\) in \(\F_p\), gives
\[
 E_p(pq+r)\equiv((p-1)!)^q r!(1+pqH_r)\pmod {p^2}.
\]
Write \(k=pj+s\). If \(2s>r\), Lucas gives \(p\mid M(pq+r,k)\),
so its square vanishes modulo \(p^2\).
Otherwise the factorial factorization gives
\[
 M(pq+r,pj+s)\equiv M(q,j)M(r,s)
 [1+p\{qH_r-2jH_s-(q-2j)H_{r-2s}\}]\pmod {p^2}.
\]
Only \(E_p\)-factors and small factorials are inverted.
The integer \(M(q,j)\) is retained exactly, so higher carries cause
no problem. Finally
\[
 (-4)^{pj+s}\equiv(-4)^{j+s}(1+pjQ_p)\pmod {p^2}.
\]
Square the multinomial congruence, multiply, and sum over \(j,s\).
\end{proof}

\begin{theorem}[Prime-square transfer]\label{transfer}
For every odd prime \(p\) and every positive integer \(m\),
\begin{equation}\label{transferformula}
                    S(p^2m)\equiv U(m)S(p^2)\pmod {p^2}.
\end{equation}
\end{theorem}
\begin{proof}
Write \(w_r=20r+17\), \(U_p=U(p)\), \(T_p=T(p)\), and \(V_p=V(p)\).
Summing \eqref{digit2} with weight \(20(pq+r)+17\) gives
\begin{equation}\label{block}
 S(pM)\equiv S(p)U(M)+pC_pT(M)+pD_pV(M)\pmod {p^2},
\end{equation}
where
\[
 C_p=20U_p+2\sum_{r<p}w_rA_r,\qquad
 D_p=4\sum_{r<p}w_rB_r .
\]
The omitted term is exactly \(pQ_p S(p)V(M)\); it vanishes because
\(p\mid S(p)\), proved in Section 2.
Lucas also gives \(L_{pu+r}\equiv a_uL_r\pmod p\), since the index
\(pv+s\) reduces to \(s\). Hence
\[
 U(pm)\equiv U_pU(m),\quad T(pm)\equiv T_pU(m),\quad
 V(pm)\equiv V_pU(m)\pmod p.
\]
Substitute \(M=pm\) into \eqref{block}. Its three coefficients
\(S(p),pC_p,pD_p\) are divisible by \(p\), so these substitutions
are valid modulo \(p^2\). The result is
\[
 S(p^2m)\equiv U(m)[S(p)U_p+pC_pT_p+pD_pV_p].
\]
The bracket is \(S(p^2)\) modulo \(p^2\) by \eqref{block} with \(M=p\).
\end{proof}

\begin{remark}[What a finite base verifies]
For a fixed prime \(p\), an exact verification of \(p^2\mid S(p^2)\)
now proves \(p^2\mid S(p^2m)\) for \emph{every} \(m\).
It is a proof for an unbounded cofactor, but not a verification over
all primes. The theorem also gives no \(p^3\) conclusion.
\end{remark}

\section{An infinite prime-square family from complex multiplication}

\begin{theorem}\label{inert}
If \(p\equiv11,13,17,19\pmod {20}\) is prime, then
\[
                     p^2\mid S(p^2m)\qquad(m\ge1).
\]
\end{theorem}
\begin{proof}
For an indeterminate \(c\), define
\[
 W_r(c)=\sum_k c^kM(r,k)^2,\qquad
 \mathcal U_p(c)=\sum_{r<p}W_r(c).
\]
Set \(h=(p-1)/2\) and \(d=\lfloor p/4\rfloor\).
The binomial-square summation in Section 2, followed by \(k=h-j\),
gives the polynomial identity
\begin{equation}\label{Upoly}
 \mathcal U_p(c)\equiv
 c^h\sum_{j=0}^{d}\frac{(4j)!}{j!^4(256c)^j}\pmod p.
\end{equation}
All apparent negative powers cancel. The finite Legendre identity
\[
 P_d(t)^2=\sum_{j=0}^d\binom dj\binom{d+j}j\binom{2j}j
                   \left(\frac{t^2-1}{4}\right)^j
\]
and \(\{-d,d+1\}=\{1/4,3/4\}\) in \(\F_p\) imply
\begin{equation}\label{square}
 \mathcal U_p(c)\equiv c^hP_d(t)^2,\qquad t^2=1-\frac1c.
\end{equation}
This is an identity in the parameter, not merely an equality at a
particular numerical value.

Take \(t_0=\sqrt5/2\), corresponding to \(c=-4\), and consider
\[
 E_{t_0}:\quad y^2=x\left(x^2+2x+\frac{1-t_0}{2}\right).
\]
Its invariants are
\[
 \Delta=8(t_0-1)^2(t_0+1),\qquad
 j(E_{t_0})=
 \frac{64(3t_0+5)^3}{(t_0-1)^2(t_0+1)}
 =632000+282880\sqrt5.
\]
Duke \cite{Duke} identifies this value as \(j(i\sqrt5)\); hence the
curve has complex multiplication of discriminant \(-20\).
The other sign of \(t_0\) gives the conjugate CM value.
The coefficients are integral away from \(2\), and
\((t_0-1)(t_0+1)=1/4\), so the model has good reduction at every
prime above an odd rational prime.

The stated residue classes are exactly
\(\left(\frac{-5}{p}\right)=-1\). Deuring's reduction theorem therefore
makes this reduction supersingular; see \cite{EOY}.
This holds over \(\overline{\F}_p\), including when \(t_0\notin\F_p\).

For completeness, the Hasse coefficient can be read off directly.
Let \(b=(1-t)/2\). The coefficient of \(x^{p-1}\) in
\([x(x^2+2x+b)]^h\) is
\[
 \sum_{k=0}^{d}\frac{h!}{k!^2(h-2k)!}\,2^{h-2k}b^k
 =
 2^h\sum_{k=0}^{d}
 \frac{(-h/2)_k((1-h)/2)_k}{k!^2}b^k
 \equiv2^hP_d(t)\pmod p.
\]
Here \((z)_k\) is the rising factorial.
The standard Hasse-invariant criterion \cite[Theorem V.4.1(a)]{Silverman}
therefore gives \(P_d(t_0)=0\). Equation \eqref{square} yields
\begin{equation}\label{Uzero} U_p=\mathcal U_p(-4)\equiv0\pmod p.\end{equation}

There are two elementary ways to obtain the other moments.
The reflection used in Section 2 also gives, for each \(k\),
\[
 \sum_{r<p}r\binom r{2k}^2
 \equiv\left(k-\frac12\right)\sum_{r<p}\binom r{2k}^2\pmod p.
\]
Thus \(T_p\equiv V_p-U_p/2\pmod p\).
Since \(20T_p+17U_p=S(p)\equiv0\pmod p\) and \(p\ne2,5\),
\eqref{Uzero} implies \(T_p=V_p=0\) in \(\F_p\).
Alternatively, differentiate \eqref{square}: the square factor gives
\(\mathcal U_p'(-4)=0\) because \(t_0\ne0\); then
\(V_p=-4\mathcal U_p'(-4)=0\).

Finally the bracket in the proof of Theorem~\ref{transfer} vanishes
modulo \(p^2\): \(p\mid S(p)\) and \(U_p,T_p,V_p\) all vanish modulo \(p\).
This proves the theorem for every cofactor \(m\).
\end{proof}

\section{An integral analytic digit operator}\label{sec:analytic}

The fixed-modulus expansion in Section~\ref{sec:transfer} is replaced here
by an exact identity in a complete ring. This keeps the entire remainder
inside a controlled integral space and permits induction through every
prime-power exponent.

Fix a prime \(p\ge5\) and \(c\in\Z_p^\times\). Write
\[
 f_c(n,k)=c^k\left(\frac{n!}{k!^2(n-2k)!}\right)^2
 \quad(0\le2k\le n),\qquad f_c(n,k)=0\quad\hbox{otherwise}.
\]
Let \(\mathcal A_p=\Z_p\langle X,Y\rangle\) denote the restricted
power-series ring: coefficients are \(p\)-adic integers and tend to zero
as total degree tends to infinity. It is complete for the Gauss norm,
and its elements may be evaluated at any pair in \(\Z_p^2\). Define
\[
 I_N(F)=\sum_{n=0}^{N-1}\sum_{0\le2k\le n}f_c(n,k)F(n,k),
 \quad U_c(N)=I_N(1),\quad T_c(N)=I_N(X),\quad V_c(N)=I_N(Y).
\]
For \(c=-4\), these agree with the earlier \(U,T,V\).

\begin{lemma}[Integral interpolation of the unit factorial]\label{an:unit}
Put
\[
 E_p(N)=\prod_{\substack{1\le i\le N\\p\nmid i}}i,\qquad W=(p-1)!.
\]
There are units \(\mathcal H_r(X)\in1+p\Z_p\langle X\rangle\),
\(0\le r<p\), such that, for every integer \(q\ge0\),
\[
 E_p(pq+r)=W^q r!\,\mathcal H_r(q),\qquad
 \mathcal H_r(X)\equiv1+pH_rX\pmod{p^2}.
\]
Here \(H_r=\sum_{i=1}^r1/i\), with \(H_0=0\). The first identity is exact.
\end{lemma}
\begin{proof}
Let \(\eta_j=\sum_{i=1}^{p-1}i^{-j}\), and set
\[
 B(X)=\prod_{i=1}^{p-1}(1+pX/i),\qquad
 \log B(X)=\sum_{j\ge1}\ell_jX^j,\quad
 \ell_j=\frac{(-1)^{j-1}p^j\eta_j}{j}.
\]
Let \(S_j(q)=\sum_{h=0}^{q-1}h^j\). Its polynomial continuation is
\[
 S_j(X)=\sum_{a=0}^j
 \left\{\begin{matrix}j\\a\end{matrix}\right\}
 \frac{(X)_{\underline{a+1}}}{a+1},
\]
where braces denote Stirling numbers of the second kind and the
underlined subscript denotes a falling factorial. Consequently
\[
 S_j(X)\in p^{-\lfloor\log_p(j+1)\rfloor}\Z_p[X].
\]
Define \(A(X)=\sum_{j\ge1}\ell_jS_j(X)\).
We claim \(A\in p^3\Z_p\langle X\rangle\).
Indeed, \(\eta_2\equiv0\pmod p\), and pairing \(i\) with \(p-i\) gives
\(2\eta_1\equiv-p\eta_2\pmod{p^2}\).
Thus the terms \(j=1,2\) have coefficient valuation at least three.
For \(j\ge3\) every coefficient of the \(j\)-th term has valuation at least
\[
 j-v_p(j)-\lfloor\log_p(j+1)\rfloor\ge3,
\]
and this bound tends to infinity. For \(j=3,4,5\) the inequality is
immediate, including the worst cases \(p=5,j=4,5\). For \(j\ge6\),
use \(v_p(j)\le\lfloor\log_5(j+1)\rfloor\) and
\(2\lfloor\log_5(j+1)\rfloor\le j-3\).
The latter is direct below \(24\); if \(t=\lfloor\log_5(j+1)\rfloor\ge2\),
then \(j\ge5^t-1\) and \(2t\le5^t-4\).

It follows that \(\mathcal E(X)=\exp A(X)\) belongs to
\(1+p^3\Z_p\langle X\rangle\): the \(k\)-th exponential term has
Gauss valuation at least \(3k-v_p(k!)\), tending to infinity.
Set
\[
 \mathcal H_r(X)=\mathcal E(X)\prod_{i=1}^r(1+pX/i).
\]
The first-order formula follows. At a nonnegative integer \(q\),
convergence permits evaluation and interchange with the finite sum:
\[
 A(q)=\sum_{h=0}^{q-1}\log B(h),\qquad
 \mathcal E(q)=\prod_{h=0}^{q-1}B(h).
\]
Splitting the unit factorial into its \(q\) complete blocks and its
last \(r\) factors proves the exact identity. We interpolate only the
normalized factor; no \(p\)-adic interpolation of \(W^q\) is required.
\end{proof}

\begin{lemma}[Exact carry-aware kernel]\label{an:kernel}
For \(0\le r,s<p\), put
\[
 d=\left\lceil\frac{2s-r}{p}\right\rceil\in\{0,1,2\},
 \qquad t=r-2s+dp,\qquad z=X-2Y.
\]
Then \(0\le t<p\). Define
\[
 \mathcal C(Y)=\exp\bigl(Y\log(c^{p-1})\bigr),\qquad
 Q_c=\frac{c^{p-1}-1}{p},
\]
and
\[
 K_{rs}(X,Y)=
 p^{2d}(z)_{\underline d}^{\,2}
 \frac{c^s W^{2d}r!^2}{s!^4t!^2}
 \frac{\mathcal C(Y)\mathcal H_r(X)^2}
 {\mathcal H_s(Y)^4\mathcal H_t(X-2Y-d)^2}.
\]
This is an element of \(\mathcal A_p\). For every \(q\ge0\) and
\(0\le2j\le q\),
\[
 f_c(pq+r,pj+s)=f_c(q,j)K_{rs}(q,j).
\]
The equality includes inadmissible original summands, interpreted as zero.
Furthermore,
\[
 K_{rs}(0,0)=f_c(r,s).
\]
If \(d\ge1\), then
\(K_{rs}\in p^2\mathcal A_p\); if \(d=0\), then
\[
 \begin{split}
 K_{rs}(X,Y)\equiv f_c(r,s)\bigl[1+p\{&
 2X(H_r-H_t)\\
 &+Y(Q_c+4H_t-4H_s)\}\bigr]\pmod{p^2}.
 \end{split}
\]
\end{lemma}
\begin{proof}
Since \(c^{p-1}\in1+p\Z_p\), the exponential defining \(\mathcal C\)
converges and has integral restricted coefficients:
\(k-v_p(k!)\to\infty\). It is congruent to \(1+pQ_cY\) modulo \(p^2\).
All factorial denominators in the kernel are units. Affine substitutions
with integral coefficients preserve restricted series, and the
\(\mathcal H\)'s are units.

Write \(q_0=q-2j\). If \(q_0\ge d\), use
\(N!=p^{\lfloor N/p\rfloor}\lfloor N/p\rfloor!E_p(N)\)
and Lemma~\ref{an:unit}. Before squaring, the factorial ratio equals
\[
 p^d(q_0)_{\underline d}\,
 \frac{q!}{j!^2q_0!}\,
 \frac{W^d r!}{s!^2t!}
 \frac{\mathcal H_r(q)}
 {\mathcal H_s(j)^2\mathcal H_t(q_0-d)}.
\]
The powers \(W^q\) cancel to the fixed factor \(W^d\).
Also \(c^{pj+s}=c^j c^s\mathcal C(j)\), proving the equality.
If \(0\le q_0<d\), the original term is zero and
\((q_0)_{\underline d}=0\), so the same equality holds. The denominator
at the negative argument \(q_0-d\) remains an analytic unit.
No \(j\) with \(2j>q\) can contribute, since
\(p(q-2j)+r-2s<0\). This proves that a common high-index range is valid
for every pair of low digits. The remaining assertions follow from
the kernel and the first-order unit expansions.
\end{proof}

\begin{proposition}[Exact analytic block identity]\label{an:block}
For \(F\in\mathcal A_p\), define
\[
 (\mathcal T_pF)(X,Y)=
 \sum_{r,s=0}^{p-1}K_{rs}(X,Y)F(pX+r,pY+s).
\]
Then \(\mathcal T_pF\in\mathcal A_p\), and
\[
 I_{pN}(F)=I_N(\mathcal T_pF)\qquad(N\ge1).
\]
There exist \(\alpha_F,\beta_F\in\Z_p\) and \(R_F\in\mathcal A_p\) such that
\begin{equation}\label{an:operator}
 \mathcal T_pF=I_p(F)+p\alpha_FX+p\beta_FY+p^2R_F.
\end{equation}
In particular, the following is an exact equality, not a truncated
congruence:
\begin{equation}\label{an:exact}
 I_{pN}(F)=I_p(F)U_c(N)+p\alpha_FT_c(N)+p\beta_FV_c(N)
                         +p^2I_N(R_F).
\end{equation}
\end{proposition}
\begin{proof}
Partition the two summation indices into their high and low digits and
apply Lemma~\ref{an:kernel}. For the expansion use
\[
 F(pX+r,pY+s)\equiv
 F(r,s)+pXF_X(r,s)+pYF_Y(r,s)\pmod{p^2}.
\]
This holds first for polynomials and then for restricted series by
Gauss-norm continuity. The kernel has an affine first-order part.
Its exact value at \((0,0)\) makes the constant term exactly \(I_p(F)\).
The residual after removing the constant and the affine first-order
part is divisible by \(p^2\) in the restricted ring. Evaluating the
resulting identity gives \eqref{an:exact}.
\end{proof}

\begin{theorem}[Simultaneous valuation induction]\label{an:moment}
Suppose \(U_c(p),T_c(p),V_c(p)\in p\Z_p\). For every \(e\ge1\), \(m\ge1\),
and \(F\in\mathcal A_p\),
\[
 I_{p^em}(F)\in p^{e-1}\Z_p.
\]
If in addition \(I_p(F)\in p\Z_p\), then
\[
 I_{p^em}(F)\in p^e\Z_p.
\]
In particular this stronger conclusion applies to \(F=1,X,Y\).
\end{theorem}
\begin{proof}
Fix \(m\). At \(e=1\), the first assertion is integrality. If
\(I_p(F)\in p\Z_p\), every term of \eqref{an:exact} with \(N=m\)
is divisible by \(p\), proving the second assertion.

Assume both assertions at \(e-1\), and set \(N=p^{e-1}m\).
By the hypothesis and the stronger induction assertion,
\(U_c(N),T_c(N),V_c(N)\in p^{e-1}\Z_p\). By the universal assertion,
\(I_N(R_F)\in p^{e-2}\Z_p\). In \eqref{an:exact}, the first term
is therefore divisible by \(p^{e-1}\) and the other three by \(p^e\).
If \(I_p(F)\) is divisible by \(p\), the first term is also divisible
by \(p^e\). This proves both assertions simultaneously, for every
restricted test function and every exponent.
\end{proof}

\begin{corollary}[All inert prime powers]\label{an:inert}
For every prime \(p\equiv11,13,17,19\pmod{20}\), every \(e\ge1\),
and every \(m\ge1\),
\[
 p^e\mid U(p^em),\qquad p^e\mid T(p^em),\qquad
 p^e\mid V(p^em),\qquad \boxed{p^e\mid S(p^em)}.
\]
\end{corollary}
\begin{proof}
Take \(c=-4\). The CM and Legendre-square argument in
Theorem~\ref{inert} proves \(U(p)\equiv0\pmod p\).
Differentiating that coefficientwise square identity gives
\(V(p)\equiv0\pmod p\); the reflection relation gives
\(T(p)\equiv V(p)-U(p)/2\equiv0\pmod p\).
Apply Theorem~\ref{an:moment}, and use \(S=20T+17U\).
\end{proof}

\section{The entire 5-part}\label{sec:five}

The prime \(5\) is handled by a different cancellation in the same
exact operator; its first three prime-block moments do not all vanish.

\begin{theorem}[One-moment criterion]\label{five:criterion}
Suppose \(p\ge5\) and
\[
 \mathcal T_p(1)=C+p^2R,\qquad C\in p\Z_p,\quad R\in\mathcal A_p.
\]
Then for all \(e,m\ge1\),
\[
 I_{p^em}(F)\in p^{e-1}\Z_p\quad(F\in\mathcal A_p),\qquad
 U_c(p^em)\in p^e\Z_p.
\]
If \(F\) is constant modulo \(p\), then \(I_{p^em}(F)\in p^e\Z_p\).
\end{theorem}
\begin{proof}
The assertions at \(e=1\) follow from integrality and
\(U_c(pm)=CU_c(m)+p^2I_m(R)\).
Assume the assertions at \(e-1\). In \eqref{an:exact}, the constant
term is divisible by \(p^{e-1}\), the two affine terms by
\(p\cdot p^{e-2}\), and the residual by \(p^2\cdot p^{e-2}\).
This proves the universal bound at \(e\). For \(F=1\), use instead
\(U_c(pN)=CU_c(N)+p^2I_N(R)\); both terms are divisible by \(p^e\).
Finally, if \(F=a+pG\), then \(I(F)=aU_c+pI(G)\), giving the last claim.
\end{proof}

\begin{corollary}\label{five:all}
For every \(e\ge1\) and \(m\ge1\), \(\boxed{5^e\mid S(5^em)}\).
\end{corollary}
\begin{proof}
Set \(p=5,c=-4\). The harmonic residues are
\((H_0,H_1,H_2,H_3,H_4)=(0,1,4,1,0)\) modulo \(5\), and \(Q_c=51\equiv1\).
For each no-carry pair put \(t=r-2s\), \(f=f_{-4}(r,s)\).
The kernel shows that the two first-order coefficients of
\(\mathcal T_5(1)\) are the sums of
\(2f(H_r-H_t)\) and \(f[Q_c+4(H_t-H_s)]\).
Their complete row-by-row evaluation is
\[
\begin{array}{c|rrr}
 r&\sum_s f&\sum_s2f(H_r-H_t)&\sum_s f[Q_c+4(H_t-H_s)]\\ \hline
0&1&0&1\\
1&1&0&0\\
2&0&2&0\\
3&2&0&1\\
4&1&3&3\\ \hline
\text{total modulo }5&0&0&0
\end{array}
\]
All carry terms are divisible by \(25\). The exact constant is
\(U(5)=1+1-15-143+1=-155\). Hence
\[
                 \mathcal T_5(1)=-155+25R,\qquad R\in\mathcal A_5.
\]
Apply Theorem~\ref{five:criterion} to \(F=20X+17\), which is constant
modulo \(5\). Notice that \(T(5)\equiv V(5)\equiv1\pmod5\);
no vanishing of those moments is being assumed.
\end{proof}

\section{Propagation through every finite chain of prime-power bases}
\label{layer:section}

The exact operator also removes a separate propagation conjecture from
the reduction to prime powers. Only its constant reduction modulo \(p\)
is needed here, so the argument includes \(p=3\).

\begin{lemma}\label{layer:odd}
For every odd prime \(p\) and \(c\in\Z_p^\times\), the exact kernel and
block identity of Lemma~\ref{an:kernel} and Proposition~\ref{an:block}
exist over \(\mathcal A_p\), with
\[
 \mathcal T_pF=I_p(F)+pR_F,\qquad R_F\in\mathcal A_p,\quad R_F(0,0)=0.
\]
At \(p=3\) no affine assertion modulo \(p^2\) is made.
\end{lemma}
\begin{proof}
For \(p\ge5\) this is already proved. For \(p=3\), the construction of
\(A(X)\) in Lemma~\ref{an:unit} still converges, now in
\(3\Z_3\langle X\rangle\). The term \(j=1\) has valuation at least two,
since \(\eta_1=3/2\). For \(j\ge2\) the Stirling bound gives
\[
 j-v_3(j)-\lfloor\log_3(j+1)\rfloor\ge1,
\]
with a bound tending to infinity. Check \(j=2,3\) directly; for \(4\le j\le7\)
the bound is at least three, and for \(j\ge8\) use
\(2\lfloor\log_3(j+1)\rfloor\le j-1\).
The exponential of an element of \(3\Z_3\langle X\rangle\) lies in
\(1+3\Z_3\langle X\rangle\), because \(k-v_3(k!)\to\infty\).
Thus all exact interpolation and carry identities remain valid,
and \(K_{rs}\equiv f_c(r,s)\pmod p\).
Finally \(F(pX+r,pY+s)\equiv F(r,s)\pmod p\), proving the assertion.
\end{proof}

\begin{theorem}[Propagation from preceding bottom layers]\label{layer:prop}
Let \(p\) be odd, \(c\in\Z_p^\times\), \(F\in\mathcal A_p\), and \(e\ge1\).
Suppose
\[
 I_{p^j}(F)\in p^j\Z_p\qquad(1\le j<e).
\]
Then, for every \(m\ge1\),
\[
 I_{p^em}(F)\equiv I_{p^e}(F)U_c(m)\pmod{p^e}.
\]
If the base divisibility also holds at \(j=e\), then
\(I_{p^em}(F)\in p^e\Z_p\) for all \(m\).
\end{theorem}
\begin{proof}
Since \(I_1(G)=G(0,0)\), repeated exact block decomposition gives
\[
                 (\mathcal T_p^jF)(0,0)=I_{p^j}(F).
\]
Inductively, if the bottom divisibilities hold through \(j-1\),
write \(\mathcal T_p^{j-1}F=p^{j-1}G\). Lemma~\ref{layer:odd} gives
\[
 \mathcal T_p^jF
 =p^{j-1}I_p(G)+p^jR_G
 =I_{p^j}(F)+p^jR_G.
\]
If the next bottom value is divisible by \(p^j\), the entire iterate
is in \(p^j\mathcal A_p\). Applying this induction through \(e-1\)
and then once more yields
\(\mathcal T_p^eF=I_{p^e}(F)+p^eR_G\).
Apply \(I_m\) to obtain the asserted congruence.
For \(e=1\) the hypothesis is empty and the same argument applies.
\end{proof}

\begin{corollary}[Reduction of the original conjecture]\label{layer:equiv}
Sun's original assertion \(n\mid S(n)\) for every \(n\ge1\) is equivalent
to its restriction \(p^e\mid S(p^e)\) to all prime powers.
\end{corollary}
\begin{proof}
Necessity is immediate. For sufficiency, use \(c=-4\) and
\(F=20X+17\) in Theorem~\ref{layer:prop} at each odd prime divisor
\(p^e\Vert n\). The \(2\)-part is Theorem~\ref{twoadic}.
The resulting prime-power divisors are pairwise coprime.
\end{proof}

The supplied manuscript \cite{Supplied} formulates an unconditional
all-layer congruence as a separate conjecture. Theorem~\ref{layer:prop}
proves precisely the propagation needed once preceding bottom layers
are known, so that separate conjecture is no longer needed for the
prime-power reduction. We do not claim its unconditional version
without those preceding hypotheses. The still missing assertions are
the unrestricted bottom values at the remaining split primes.

\section{Catalan coordinates remove the off-diagonal condition}
\label{cat:section}

The negative-index continuation in the supplied manuscript \cite{Supplied} introduces a
two-variable divisibility condition. Its two stated one-variable conditions
are not independent: the diagonal condition alone implies the off-diagonal
one. We prove this by an integral change of variables. The diagonal condition
itself remains unproved in full.

Put
\[
 \gamma_l=(-4)^l\binom{2l}{l}^{\!2},\qquad
 G(x)=\sum_{l\ge0}\gamma_lx^{2l},\qquad
 \rho=\frac{uv}{1+u+v},\qquad
 \Lambda=\frac{(1+u)(1+v)}{uv},
\]
and define
\[
 A=\Z[[u,v]],\qquad
 \theta_u=u\frac{\partial}{\partial u},\quad
 \theta_v=v\frac{\partial}{\partial v},\qquad
 \mathcal D_A=\theta_u A+\theta_v A,
 \qquad \Psi=(3\rho+20\rho^2)G(uv).
\]
Constant terms are taken in the ring
\(\bigcup_{N\ge0}(uv)^{-N}A\), so every operation below is coefficientwise
well-defined.

\begin{lemma}[The negative-index reduction]\label{cat:negative}
For every positive integer \(N\),
\begin{equation}\label{cat:negativeidentity}
 S(N)=-\CT(\Lambda^N\Psi)
       +20N\,\CT\bigl(\Lambda^N\rho G(uv)\bigr).
\end{equation}
Consequently, \(\Psi\in\mathcal D_A\) implies \(N\mid S(N)\) for every
\(N\ge1\).
\end{lemma}
\begin{proof}
First,
\[
 \CT\bigl(\Lambda^kG(uv)\bigr)
 =\sum_l\gamma_l\binom{k}{2l}^{\!2}=a_k.
\]
The polynomial identity
\[
 (T-1)^2\sum_{k=0}^{N-1}(20k+17)T^k
 =(3T+17)(1-T^N)+20NT^N(T-1)
\]
may be evaluated at \(T=\Lambda\), since \((\Lambda-1)^{-1}=\rho\).
Multiplication by \(G(uv)\) and taking constant terms give
\eqref{cat:negativeidentity}; the constant term of \(\Psi\) is zero.
Equivalently, the negative-index generating series is
\[
 \sum_{n\ge1}(20n-17)\Lambda^{-n}=3\rho+20\rho^2.
\]
If \(\Psi=\theta_uJ_1+\theta_vJ_2\) with integral series \(J_i\), then
formal integration by parts gives
\[
 \CT(\Lambda^N\theta_uJ_1)
 =-N\,\CT\bigl(\Lambda^{N-1}\theta_u(\Lambda)J_1\bigr)\in N\Z,
\]
and the same holds for the \(v\)-derivative. This proves the implication.
\end{proof}

For a series \(F\in A\), write
\[
 \operatorname{Diag}F(x)=\sum_{n\ge0}[u^nv^n]F\,x^n.
\]
We use the elementary coefficient criterion
\begin{equation}\label{cat:bezout}
 F\in\mathcal D_A
 \quad\Longleftrightarrow\quad
 [u^0v^0]F=0\ \text{ and }\
 \gcd(r,s)\mid[u^rv^s]F\quad((r,s)\ne(0,0)).
\end{equation}
Indeed, each coefficient of a divergence has the form \(rj_1+sj_2\).
Conversely, Bezout's identity permits such integers \(j_1,j_2\) to be
chosen separately for every coefficient.

\begin{theorem}[Integral diagonal criterion]\label{cat:criterion}
Let \(R(x)\in\Z[[x]]\) and \(\alpha,\beta\in\Z\). For
\[
 F(u,v)=(\alpha\rho+\beta\rho^2)R(uv),
\]
the following conditions are equivalent:
\begin{enumerate}
 \item \(F\in\mathcal D_A\);
 \item \(n\mid[u^nv^n]F\) for every \(n\ge1\).
\end{enumerate}
The statement holds over \(\Z\), including at the prime \(2\).
\end{theorem}
\begin{proof}
Only the reverse implication requires proof. Let \(C(x)\in1+x\Z[[x]]\)
be the Catalan series, characterized by \(C=1+xC^2\). Define
\[
 U=uC(uv),\qquad V=vC(uv),\qquad y=UV.
\]
This is an integral formal automorphism with inverse
\begin{equation}\label{cat:inverse}
 u=\frac{U}{1+y},\qquad v=\frac{V}{1+y},\qquad
 x=uv=\frac{y}{(1+y)^2},\qquad C(x)=1+y.
\end{equation}
In these coordinates,
\[
 1+u+v=\frac{(1+U)(1+V)}{1+y},\qquad
 \rho=\frac{y}{(1+y)(1+U)(1+V)}.
\]
The logarithmic differential matrix of the inverse is
\[
 \begin{pmatrix}d\log u\\d\log v\end{pmatrix}
 =\frac1{1+y}
 \begin{pmatrix}1&-y\\-y&1\end{pmatrix}
 \begin{pmatrix}d\log U\\d\log V\end{pmatrix}.
\]
Its determinant \(J(y)=(1-y)/(1+y)\) is an integral power-series unit;
the matrix and its inverse have integral coefficients. Exact logarithmic
two-forms are therefore preserved in both directions. Explicitly,
\[
 F\,d\log u\wedge d\log v
 =d\bigl(J_1\,d\log v-J_2\,d\log u\bigr)
\]
is equivalent to \(F=\theta_uJ_1+\theta_vJ_2\), and pulling back the
one-form keeps both coefficients integral. Hence \(F\in\mathcal D_A\)
if and only if
\[
 \widetilde F(U,V)
 =J(y)F\left(\frac{U}{1+y},\frac{V}{1+y}\right)
\]
is an integral formal divergence in \(U,V\).

There are integral series
\[
 A_1(y)=\frac{\alpha y(1-y)}{(1+y)^2}
          R\left(\frac{y}{(1+y)^2}\right),\qquad
 A_2(y)=\frac{\beta y^2(1-y)}{(1+y)^3}
          R\left(\frac{y}{(1+y)^2}\right)
\]
such that
\begin{equation}\label{cat:separated}
 \widetilde F(U,V)=
 \frac{A_1(y)}{(1+U)(1+V)}
 +\frac{A_2(y)}{(1+U)^2(1+V)^2}.
\end{equation}
Write \(A_1(y)=\sum_k A_{1,k} y^k\), \(A_2(y)=\sum_k A_{2,k} y^k\).
For \(r,s\ge1\), put \(q=\min(r,s)\) and \(d=\gcd(r,s)\).
Expanding the separated denominators gives
\[
 [U^rV^s]\widetilde F
 =(-1)^{r+s}\sum_{k=0}^q
       \bigl(A_{1,k}+A_{2,k}(r-k+1)(s-k+1)\bigr).
\]
Since \(r,s,q\) are all divisible by \(d\),
\begin{equation}\label{cat:offdiagonal}
 [U^rV^s]\widetilde F
 \equiv(-1)^{r+s}[U^qV^q]\widetilde F\pmod d.
\end{equation}
Thus divisibility of the diagonal coefficients is sufficient for all
the required off-diagonal divisibilities in these coordinates.

It remains to justify diagonal divisibility after the change of variables.
The radial substitution \eqref{cat:inverse} preserves the difference of
the two exponents. Therefore only old diagonal monomials contribute to
the new diagonal, and
\[
 \operatorname{Diag}\widetilde F(y)
 =J(y)\,\operatorname{Diag}F\left(\frac{y}{(1+y)^2}\right).
\]
By the assumed diagonal divisibility there is \(H(x)\in x\Z[[x]]\)
with \(\operatorname{Diag}F(x)=xH'(x)\). Since
\[
 J(y)=\frac{y}{x(y)}\frac{dx(y)}{dy},\qquad
 x(y)=\frac{y}{(1+y)^2},
\]
we obtain
\[
 \operatorname{Diag}\widetilde F(y)
 =y\frac{d}{dy}H(x(y)).
\]
The composition \(H(x(y))\) is integral, so its displayed derivative
has its \(q\)-th coefficient divisible by \(q\). Equation
\eqref{cat:offdiagonal} proves all gcd divisibilities. The constant and
axis coefficients vanish because \(A_1\in y\Z[[y]]\) and
\(A_2\in y^2\Z[[y]]\). Apply \eqref{cat:bezout}, and then the integral
coordinate invariance, to conclude \(F\in\mathcal D_A\).
\end{proof}

\begin{corollary}[A single sufficient condition]\label{cat:single}
For \(n\ge1\), put
\[
 Y(n)=\sum_{l=0}^{\lfloor(n-1)/2\rfloor}
  \gamma_l\binom{2i}{i}(10i+3),\qquad i=n-1-2l.
\]
Then
\[
 \Psi\in\mathcal D_A
 \quad\Longleftrightarrow\quad n\mid Y(n)\quad(n\ge1).
\]
In particular, the latter statement alone implies the original assertion
\(n\mid S(n)\) for every \(n\).
\end{corollary}
\begin{proof}
The diagonal coefficient of \(\Psi\) is exactly \(Y(n)\). Indeed, the
contribution with fixed \(l\), and \(i=n-1-2l\), is
\[
 \gamma_l\left(3\binom{2i}{i}
       +20(2i-1)\binom{2i-2}{i-1}\right)
 =\gamma_l(10i+3)\binom{2i}{i}.
\]
For \(i=0\) the second summand is absent. For \(i\ge1\), the equality
uses \(20(2i-1)\binom{2i-2}{i-1}=10i\binom{2i}{i}\).
Apply Theorem~\ref{cat:criterion} with
\((\alpha,\beta,R)=(3,20,G)\), then Lemma~\ref{cat:negative}.
\end{proof}

The remaining single condition can also be written
\begin{equation}\label{cat:primitive}
 \int_0^x\frac{3+8t}{(1-4t)^{3/2}}G(t)\,dt\in\Z[[x]],
\end{equation}
because the integrand is \(\sum_{n\ge1}Y(n)x^{n-1}\). The integral is a
formal antiderivative over \(\Q[[x]]\); its integrality is precisely the
assertion to be proved, not an assumption silently used in the argument.

\begin{corollary}[The off-diagonal condition is redundant]\label{cat:redundant}
Assume \(n\mid Y(n)\) for every \(n\ge1\). For every \(1\le b<a\),
\[
 \gcd(a,b)\ \Bigm|\
 \sum_{l=0}^{\lfloor(b-1)/2\rfloor}
 \gamma_l\binom{2(a-1-2l)}{a-1-2l}
                  \bigl(10(a-1-2l)+3\bigr).
\]
Thus condition (II) of the supplied manuscript follows from its
condition (I).
\end{corollary}
\begin{proof}
By Corollary~\ref{cat:single}, \(\Psi\) is an integral formal divergence.
Make the integral coordinate change
\[
 \varphi=u(1+u),\qquad \chi=\frac{v}{1+u},\qquad
 E(\varphi)=(1+4\varphi)^{-1/2}.
\]
The inverse of \(u\mapsto u(1+u)\) is an integral Catalan series,
and the logarithmic Jacobian \((1+2u)/(1+u)\) is a unit.
Direct substitution shows that the same logarithmic two-form has
coefficient
\[
 \Theta(\varphi,\chi)=
 \varphi\chi\,
 \frac{10\chi+(3-7\chi)E(\varphi)}{(1+\chi)^2}
 G(\varphi\chi)
\]
in the new coordinates. Hence \(\Theta\) is also an integral formal
divergence. For \(b<a\), all the indices \(i=a-1-2l\) occurring below
are positive. Expanding
\(E(\varphi)=\sum_i(-1)^i\binom{2i}{i}\varphi^i\) gives exactly
\[
 [\varphi^a\chi^b]\Theta
 =(-1)^{a+b}\sum_{l=0}^{\lfloor(b-1)/2\rfloor}
       \gamma_l\binom{2i}{i}(10j+3),
 \qquad i=a-1-2l,\quad j=b-1-2l.
\]
Since \(j-i=b-a\) is divisible by \(\gcd(a,b)\), replacing \(j\) by
\(i\) does not change this expression modulo that gcd. The coefficient
itself is divisible by the gcd, by \eqref{cat:bezout}. This proves the
claimed congruence.
\end{proof}

The result removes an auxiliary condition; it does not prove the
remaining integrality statement \eqref{cat:primitive}. We have proved
that this statement is equivalent to the two-variable divergence
criterion and is sufficient for Sun's original conjecture. We do not
assert the converse implication from the original conjecture to
\eqref{cat:primitive}.
\section{A complete unweighted branch: a quadratic invariant}\label{zeroliftsec}

The following statement is separate from the weighted all-\(n\) target.
It supplies a full proof of the zero branch of Sun's Conjecture 5.8(i)
\cite{Sun2010}, recorded as \texttt{OeisA179524.conjecture3} in
\cite{Formal}. Its historical priority is not asserted.

\begin{theorem}[Automatic lift away from singular parameters]\label{zerolift}
Let \(p\) be an odd prime and set
\(Y_p(c)=\sum_{r=0}^{p-1}W_r(c)\). If \(c_0\in\Z_p\),
\(c_0(c_0-1)\in\Z_p^\times\), and \(Y_p(c_0)\equiv0\pmod p\), then
\[
                         Y_p(c_0)\equiv0\pmod {p^2}.
\]
\end{theorem}

The binomial reductions used below are established ingredients:
Sun \cite[Theorem 1.1]{SunApery} relates this sum to a sum of Ap\'ery
polynomials, and Guo--Zeng \cite[Theorem 1.3]{GuoZeng} give the
binomial reduction used here. We include its direct proof to make the
subsequent differential argument checkable without additional files.
The finite-invariant lifting strategy itself has a direct antecedent in
Babanskyy \cite[Lemmas 2.1--2.2]{Babanskyy}: a quadratic differential
invariant, explicit treatment of exceptional degree-\(p\) coefficients,
and a double-root argument yield a lift from modulo \(p\) to modulo
\(p^2\). Here that strategy is adapted to the dual polynomial below.
The proof is given directly; the invariant principle is not claimed
as a new invention.

\begin{lemma}\label{vandermonde}
For an even integer \(m\) with \(0\le m\le p-1\),
\[
 \sum_{r=m}^{p-1}\binom rm^2
                \equiv\binom{p+m}{2m+1}\pmod {p^2}.
\]
\end{lemma}
\begin{proof}
For \(m\le r<p\), all denominators in \(H_r-H_{r-m}\) are \(p\)-adic
units, and
\[
 \binom{p+m-1-r}{m}\equiv
 \binom rm[1-p(H_r-H_{r-m})]\pmod {p^2}.
\]
The sign is positive because \(m\) is even.
The involution \(r\mapsto p+m-1-r\) permutes \(m,\ldots,p-1\),
preserves \(\binom rm^2\) modulo \(p\), and reverses
\(H_r-H_{r-m}\) modulo \(p\). The last fact follows from
\(H_{p-1-u}\equiv H_u\pmod p\). Thus the sum of the first-order
corrections is zero modulo \(p\).
Vandermonde's identity gives
\[
 \sum_{r=m}^{p-1}\binom rm\binom{p+m-1-r}{m}
                  =\binom{p+m}{2m+1},
\]
proving the claim, including \(m=0\) and \(m=p-1\).
\end{proof}

\begin{lemma}[Dual polynomial and differential equation]\label{dual}
Put \(h=(p-1)/2\) and
\[
 Z_p(c)=\sum_{s=0}^{h}v_s c^s,\qquad
 v_s=p\frac{(2s)!^4}{s!^4(4s+1)!}.
\]
Then \(v_s\in\Z_p\), \(Y_p\equiv Z_p\pmod {p^2}\) coefficientwise,
and, with \(\theta=c\,d/dc\),
\begin{equation}\label{ode}
 \mathcal L Y_p\equiv0\pmod {p^2},\qquad
 \mathcal L=\theta(4\theta-1)(4\theta+1)-2c(2\theta+1)^3 .
\end{equation}
\end{lemma}
\begin{proof}
The denominator \((4s+1)!\) contains at most one factor \(p\),
since \(4s+1\le2p-1\); it is canceled by the displayed factor \(p\).
Lemma~\ref{vandermonde} makes the coefficient of \(c^s\) in \(Y_p\)
congruent to
\[
 \binom{2s}s^2\binom{p+2s}{4s+1}
 =
 p\binom{2s}s^2
       \frac{\prod_{j=1}^{2s}(p^2-j^2)}{(4s+1)!}.
\]
The product is \((2s)!^2\) modulo \(p^2\), since \(2s\) is even.
Before division the error has valuation at least three; the factorial
denominator loses at most one. Hence the error is still zero modulo
\(p^2\), proving the coefficientwise congruence.

The exact coefficient relation is
\[
 s(4s-1)(4s+1)v_s=2(2s-1)^3v_{s-1}\qquad(1\le s\le h).
\]
Consequently the entire truncation residual is
\[
             \mathcal L Z_p=-2p^3v_h c^{h+1}.
\]
It vanishes modulo \(p^3\). Applying the integer differential operator
to \(Y_p-Z_p\) proves \eqref{ode}. No coefficient \(s\) or \(4s\pm1\)
has been inverted modulo \(p\).
\end{proof}

\begin{lemma}[Quadratic differential invariant]\label{invariant}
For a polynomial \(Y\), define
\[
 \mathcal J[Y]=16c^2(c-1)(Y'^2-2YY'')
              -16c(3c-2)YY'-(4c+1)Y^2.
\]
Then
\[
             \frac{d}{dc}\mathcal J[Y]=2Y\,\frac{\mathcal LY}{c}.
\]
In particular, \(\mathcal J[Y_p]\equiv0\pmod {p^2}\) coefficientwise.
\end{lemma}
\begin{proof}
The differential identity follows by expanding
\[
 \frac{\mathcal LY}{c}
 =16c^2(1-c)Y'''+(48c-72c^2)Y''
                  +(15-52c)Y'-2Y.
\]
The right side is polynomial. Thus \(\mathcal J[Y_p]'\) is zero
modulo \(p^2\). A zero derivative alone would not suffice, because a
degree-\(p\) term could survive. We check that term explicitly.

Since \(\deg Y_p=h\), the degree of \(\mathcal J[Y_p]\) is at most
\(2h+1=p\). Its constant coefficient is \(-Y_p(0)^2=-p^2\).
If \(K\) is the leading coefficient of \(Y_p\), its degree-\(p\)
coefficient is
\[
 [16(h^2-2h(h-1))-48h-4]K^2
                   =-4(2h+1)^2K^2=-4p^2K^2.
\]
Both endpoints vanish modulo \(p^2\). For every intermediate index
\(1\le j\le p-1\), the derivative gives \(j\, [c^j]\mathcal J[Y_p]=0\)
modulo \(p^2\), and \(j\) is a unit. All coefficients therefore vanish.
\end{proof}

\begin{proof}[Proof of Theorem~\ref{zerolift}]
Equation \eqref{square} gives
\[
 Y_p(c)\equiv c^hP_d(t)^2\pmod p,\qquad
 d=\lfloor p/4\rfloor,\quad t^2=1-1/c.
\]
Choose \(t_0\) over the algebraic closure of the residue field with
\(t_0^2=1-1/c_0\).
The assumptions make \(t_0\ne0\), \(t_0^2\ne1\), and \(P_d(t_0)=0\).
Such a Legendre root is simple. Indeed, if its multiplicity were
\(m\ge2\), the lowest-degree term of
\[
        (1-t^2)P_d''-2tP_d'+d(d+1)P_d=0
\]
would have a nonzero coefficient proportional to
\((1-t_0^2)m(m-1)\), since \(m\le d<p\).
The degree-zero case has no roots and is vacuous.

Because \(dt/dc=1/(2tc^2)\), differentiating the polynomial congruence
gives
\[
 Y_p'(c_0)\equiv0,\qquad
 Y_p''(c_0)\equiv
 2c_0^hP_d'(t_0)^2(2t_0c_0^2)^{-2}\ne0\pmod p.
\]
The expression is in the base residue field, as is the left side;
the field extension is used only to describe the root.

Evaluate the invariant at \(c_0\). Both \(Y_p(c_0)\) and
\(Y_p'(c_0)\) are divisible by \(p\), so every term except the
\(Y_pY_p''\) term vanishes modulo \(p^2\). We obtain
\[
       32c_0^2(c_0-1)Y_p(c_0)Y_p''(c_0)\equiv0\pmod {p^2}.
\]
Its coefficient multiplying \(Y_p(c_0)\) is a \(p\)-adic unit.
This forces \(Y_p(c_0)\equiv0\pmod {p^2}\).
\end{proof}

\begin{corollary}[Sun's inert unweighted branch]\label{unweighted}
For every prime \(p\equiv11,13,17,19\pmod {20}\),
\[
                 \sum_{r=0}^{p-1}a_r\equiv0\pmod {p^2}.
\]
\end{corollary}
\begin{proof}
The CM argument in Theorem~\ref{inert} gives \(Y_p(-4)=0\pmod p\).
Here \(c_0=-4\) satisfies \(c_0(c_0-1)=20\), a unit at all the
specified primes. Apply Theorem~\ref{zerolift}.
\end{proof}

The corollary closes all quantifiers of this particular unweighted
branch. It neither proves the split-prime evaluation branches of
Sun's Conjecture 5.8(i), nor the weighted all-\(n\) statement
\eqref{target} in Conjecture 5.8(ii).


\section{A direct symbolic proof modulo 9}

\begin{proposition}\label{nine}
For every \(m\ge1\), \(9\mid S(9m)\).
\end{proposition}
\begin{proof}
The digit expansion \eqref{digit2} at \(p=3\), or a direct factorial
expansion, gives for every \(q\ge0\)
\begin{equation}\label{ternary}
 a_{3q}\equiv a_q+6L_q,\qquad
 a_{3q+1}\equiv a_q,\qquad
 a_{3q+2}\equiv-15a_q+3L_q\pmod9.
\end{equation}
Here is an explicit check of the coefficients. The only low-digit
pairs with \(2s\le r\) and their corresponding weights are
\[
 (r,s)=(0,0),(1,0),(2,0),(2,1),\qquad f(r,s)=1,1,1,-16.
\]
Use \(H_0=0,H_1=1,H_2=3/2\), together with
\[
 (-4)^{3j+s}\equiv(-4)^{j+s}(1+6j)\pmod9.
\]
The correction divided by \(3\) is
\[
 2q(H_r-H_{r-2s})+4j(H_{r-2s}-H_s)+2j .
\]
Its \(q\)-coefficient is zero modulo \(3\) in all four cases, and its
\(j\)-coefficients are \(2,0,2,1\), respectively.
All low-carry terms vanish modulo \(9\); higher carries remain in
\(M(q,j)\). This proves \eqref{ternary} for arbitrary \(q\).

Consequently
\[
 \sum_{r=0}^2(20(3q+r)+17)a_{3q+r}
 \equiv3(q a_q+L_q)\pmod9.
\]
Modulo \(3\), Lucas gives \(a_{3u+r}\equiv a_u a_r\) and
\(L_{3u+r}\equiv a_u L_r\). Since
\[
 (a_0,a_1,a_2)\equiv(1,1,0),\qquad
 (L_0,L_1,L_2)\equiv(0,0,2)\pmod3,
\]
we have
\[
 \sum_{r=0}^2((3u+r)a_{3u+r}+L_{3u+r})
 \equiv a_u\sum_{r=0}^2(r a_r+L_r)=0\pmod3.
\]
Sum the preceding modulo-\(9\) block over \(q<3m\), and group the
right side into these three-term blocks.
\end{proof}


\section{Certified computation and the remaining bottom layers}\label{sec:computations}

Appendix~\ref{rec:section} proves an exact telescoping certificate for the
order-four recurrence stated in the supplied manuscript.
The leading recurrence coefficient never vanishes for \(n\ge0\), so
the initial values determine the original binomial-sum sequence uniquely.
Appendix~\ref{app:verify} then computes this sequence with exact integer
arithmetic, checking each division. The fresh run verifies
\(n\mid S(n)\) for every \(1\le n\le50000\), with no failure.

Every odd prime power \(p^e\le50000\) comes with its complete chain of
smaller bases. Theorem~\ref{layer:prop} therefore turns each such base
into \(p^e\mid S(p^em)\) for all \(m\ge1\).
Together with Corollaries~\ref{an:inert}, \ref{five:all},
Theorem~\ref{twoadic}, and the prime case, this proves
Theorem~\ref{main}.

Some fresh exact checks of the analytic conclusions are recorded below.
They are checks on the implementation and are not assumptions of the
unbounded-exponent proofs. Put \(M_2(N)=\sum_{n<N}n^2a_n\).
\[
\begin{array}{c|r|rrrr}
p&e\ \text{with }N=p^e&v_p(U(N))&v_p(T(N))&v_p(S(N))&v_p(M_2(N))\\ \hline
13&3&4&3&3&2\\
17&3&4&4&4&2\\
19&3&4&3&3&2\\
11&4&4&5&4&4\\
5&6&6&6&6&6\\
13&4&4&5&4&4
\end{array}
\]

For comparison with earlier formulations, define the Newton coefficients
\[
 b_\ell=\sum_{j=0}^{\ell}(-1)^{\ell-j}\binom{\ell}{j}a_j,\qquad b_{-1}=0.
\]
Binomial inversion and \eqref{hockey} give
\[
 \frac{S(n)}n=\sum_{\ell<n}\binom{n-1}{\ell}
 \frac{(20\ell+17)b_\ell+20\ell b_{\ell-1}}{\ell+1}.
\]
Since the transform and its inverse are integral, the original conjecture
is equivalent to
\[
                 \ell+1\mid3b_\ell+20b_{\ell-1}\qquad(\ell\ge0).
\]
This remains an equivalent unsolved condition. The single sufficient
integrality statement in \eqref{cat:primitive} is also not proved.

The remaining original cases are bottom values \(p^e\mid S(p^e)\)
at split primes \(p\equiv1,3,7,9\pmod{20}\), beyond their displayed
finite exponent bounds. There is now no missing all-cofactor propagation
step once a finite chain of bottom values is supplied.
Neither the two nonzero unweighted evaluations in Sun 5.8(i) nor
the separate weighted prime-square refinement
\[
 S(p)\stackrel{?}{\equiv}
 p\left(7+10\left(\frac{-1}{p}\right)\right)\pmod{p^2}
\]
is proved in general here.

\section{Provenance, review, and scope}

The exact original target is the statement recorded as
\texttt{OeisA179524.conjecture4} in the inspected formal-conjectures file
\cite{Formal}. Its first path commit is dated 20 August 2026, rather than
the November date given in the supplied manuscript. Although the
OEIS Open article \cite{OEISOpen} does not mention this sequence by its
identifier, the released benchmark includes the target. These metadata
observations are not a certificate that no later solution exists.

The supplied manuscript \cite{Supplied} establishes useful exact
reductions but states its order-four recurrence as empirically found.
Appendix~\ref{rec:section} supplies the missing all-index proof.
Its reported million-index run and prime cutoff \(6000\) are not inherited
here: the cited external computation files were not among the attachments.
The independently reproduced cutoff in this note is \(50000\), and its
complete program is embedded in the single source.

The use of analytic factorial weights and a simultaneous induction
with a one-power loss for arbitrary test functions has a substantial
antecedent in Delaygue \cite[Theorem 1, Proposition 1, and
Sections 3.3--4.3]{Delaygue}. Here we give an explicit restricted-series
operator for the triangular truncations in Sun's problem, an affine
first-order refinement, full carry control, and moment criteria implying
the stated prime-power divisibilities. No claim is made to have invented
the general analytic-weight strategy. The finite differential-invariant
method used for the separate unweighted branch also has the explicit
antecedent cited in Section~\ref{zeroliftsec}.

Independent mathematical derivations checked the exact interpolation,
the Gauss-norm bounds, every carry boundary, the simultaneous induction,
the \(p=5\) coefficient table, all-layer propagation, the Catalan change of
variables, and the recurrence certificate. This is internal mathematical
review, not external peer review or a Lean formalization.
A bounded primary-source search did not identify an exact earlier
resolution of the new target-specific conclusions; it does not establish
historical nonduplication. Thus the full original all-\(n\) success
condition remains unmet, and historical priority remains unasserted.

\appendix
\section{An exact certificate for the order-four recurrence}
\label{rec:section}

This appendix supplies a complete telescoping certificate for the
order-four recurrence stated in the supplied manuscript.  In particular,
the recurrence is established for the original binomial sum at every
nonnegative index; it is not assumed on the basis of numerical agreement.
Write
\begin{equation}\label{rec:summand}
 f(n,k)=
 \begin{cases}
 \displaystyle(-4)^k\frac{(n!)^2}{(k!)^4((n-2k)!)^2},&0\leq 2k\leq n,\\[4pt]
 0,&2k>n,
 \end{cases}
 \qquad a(n)=\sum_{k\geq0}f(n,k).
\end{equation}
All indices in this appendix are nonnegative integers unless stated
otherwise.

\subsection{The recurrence and its certificate}

Define the following five polynomials:
\begin{align}
 p_0(n)&=-4225(n+1)^2(n+2)^2
       (240n^3+2160n^2+6483n+6490),\label{rec:p0}\\
 p_1(n)&=-2(n+2)^2\bigl(91680n^5+1146000n^4+5647266n^3\notag\\
       &\hspace{38mm}{}+13701286n^2+16361708n+7698385\bigr),\label{rec:p1}\\
 p_2(n)&=-2\bigl(16080n^7+289440n^6+2219001n^5+9390796n^4\notag\\
       &\hspace{17mm}{}+23689644n^3+35620787n^2+29562269n
       +10448048\bigr),\label{rec:p2}\\
 p_3(n)&=2(n+3)\bigl(480n^6+7920n^5+53646n^4+190760n^3\notag\\
       &\hspace{34mm}{}+375516n^2+388401n+165272\bigr),\label{rec:p3}\\
 p_4(n)&=-(n+3)(n+4)^3(240n^3+1440n^2+2883n+1927).
       \label{rec:p4}
\end{align}
For a compact certificate, put
\begin{align}
 A_n&=240n^3+2160n^2+6483n+6490,\notag\\
 B_n&=240n^3+1440n^2+2883n+1927,\notag\\
 C_n&=720n^5+8640n^4+48369n^3+163272n^2+315171n+257123,
       \label{rec:abc}\\
 P_n(t)&=65A_nt^4+4(n-31)A_nt^3+2C_nt^2
          +(n+3)^3B_n(4t+n+4),\label{rec:P}\\
 G(n,k)&=\frac{16k^4}{(n+3)^2(n+4)^2}
            P_n(n+4-2k)f(n+4,k).\label{rec:G}
\end{align}
The summand identity is
\begin{equation}\label{rec:telescoper}
 \sum_{i=0}^{4}p_i(n)f(n+i,k)=G(n,k+1)-G(n,k).
\end{equation}
Consequently, the original sequence satisfies
\begin{equation}\label{rec:recurrence}
 \boxed{\displaystyle\sum_{i=0}^{4}p_i(n)a(n+i)=0
        \qquad(n\geq0).}
\end{equation}

\subsection{Polynomial verification and the endpoints}

Here is a verification of \eqref{rec:telescoper} that includes the
truncated summation boundaries.  Set
\begin{align}
 D_n&=\prod_{j=1}^{4}(n+j)^2,\notag\\
 Q(n,k)&=\sum_{i=0}^{4}p_i(n)
       \prod_{j=1}^{i}(n+j)^2
       \prod_{j=i+1}^{4}(n+j-2k)^2,\label{rec:Q}\\
 R(n,k)&=16(n+1)^2(n+2)^2P_n(n+4-2k).\label{rec:R}
\end{align}
Empty products are interpreted as~1.  Substitution of
\eqref{rec:p0}--\eqref{rec:p4} and \eqref{rec:abc}--\eqref{rec:P}
gives the polynomial identity in $\mathbb Z[n,k]$
\begin{equation}\label{rec:polynomial-identity}
 \boxed{\begin{aligned}
 Q(n,k)={}&-4(n+4-2k)^2(n+3-2k)^2R(n,k+1)\\
          &{}-k^4R(n,k).
 \end{aligned}}
\end{equation}
Thus the entire certificate can be checked by finite polynomial
multiplication and coefficient comparison.  This is an identity of
polynomials, rather than a test at finitely many numerical pairs.

For $0\leq2k\leq n+4$ one has
\begin{equation}\label{rec:common-denominator}
 D_nf(n+i,k)=f(n+4,k)
       \prod_{j=1}^{i}(n+j)^2
       \prod_{j=i+1}^{4}(n+j-2k)^2.
\end{equation}
When $2k\leq n+i$, this follows directly by cancelling factorials.
When $n+i<2k\leq n+4$, the left-hand side is zero by definition;
the right-hand side is also zero because its second product contains
the factor with $j=2k-n$.  This proves
\eqref{rec:common-denominator} even when some of the five shifted
summands have left their support.

The remaining ratio is
\begin{equation}\label{rec:ratio}
 f(n+4,k+1)=
 -4\frac{(n+4-2k)^2(n+3-2k)^2}{(k+1)^4}\,f(n+4,k).
\end{equation}
It holds throughout the support of $f(n+4,k)$, including its last
nonzero term: at that term either $n+4-2k=0$ or $n+3-2k=0$,
and both sides of \eqref{rec:ratio} are zero.  Multiplying
\eqref{rec:polynomial-identity} by $f(n+4,k)/D_n$, and using
\eqref{rec:common-denominator} and \eqref{rec:ratio}, proves
\eqref{rec:telescoper} for $0\leq2k\leq n+4$.  For $2k>n+4$
all its terms vanish.  Finally, $G(n,0)=0$ because of the factor
$k^4$, and $G(n,k)=0$ for sufficiently large $k$ by finite
support.  Summing \eqref{rec:telescoper} over $k\geq0$ therefore
proves \eqref{rec:recurrence} with no boundary remainder.

\subsection{What the certificate does and does not establish}

For every $n\geq0$ the leading coefficient $p_4(n)$ is nonzero.
Together with
\begin{equation}\label{rec:initial-values}
 (a(0),a(1),a(2),a(3))=(1,1,-15,-143),
\end{equation}
the recurrence therefore determines the original binomial-sum
sequence uniquely.  A subsequent calculation using this recurrence
and exact integer arithmetic consequently computes that original
sequence, provided the implementation and the stated range are
independently checked.

The certificate alone does not authenticate the supplied manuscript's
claim of a check through one million indices.  The computation itself
has not been recovered from that claim, and the external files
\texttt{ctest.c} and \texttt{verify\_log.txt} mentioned there were not
included in the supplied \TeX{} and PDF.  Any finite-range verification
reported in the present paper must therefore be supported by a fresh
computation or by an independently checked log.  Neither a recurrence
certificate nor a finite computation proves $N\mid S(N)$ for all~$N$.

\section{Complete exact verification program}\label{app:verify}

The following Python 3.11-or-later program uses only the standard library.
The recurrence is justified by Appendix~\ref{rec:section}; comparison
with the original definition through index \(100\) checks its implementation.
All \(49996\) recurrence divisions and \(50000\) divisibility tests passed in the
fresh run. The additional checkpoint valuations and exact density
comparison are reproduced by the same program.
The analytic all-exponent theorems do not depend on this computation.

\begin{lstlisting}
from pathlib import Path
from math import comb, isqrt
import sys, time, json
if hasattr(sys, "set_int_max_str_digits"):
    sys.set_int_max_str_digits(0)
N = 50000
def ps(n):
    return (
    -4225*(n+1)**2*(n+2)**2*(240*n**3+2160*n**2+6483*n+6490),
    -2*(n+2)**2*(91680*n**5+1146000*n**4+5647266*n**3+13701286*n**2+16361708*n+7698385),
    -2*(16080*n**7+289440*n**6+2219001*n**5+9390796*n**4+23689644*n**3+35620787*n**2+29562269*n+10448048),
    2*(n+3)*(480*n**6+7920*n**5+53646*n**4+190760*n**3+375516*n**2+388401*n+165272),
    -(n+3)*(n+4)**3*(240*n**3+1440*n**2+2883*n+1927))
def direct(n):
    return sum((-4)**k*comb(n,k)**2*comb(n-k,k)**2 for k in range(n//2+1))
def val(x,p):
    if not x:return None
    v=0
    while x%p==0:x//=p;v+=1
    return v
win=[1,1,-15,-143]
S=U=T=M2=0
checkpoints={15625:(5,6),14641:(11,4),28561:(13,4),2197:(13,3),4913:(17,3),6859:(19,3)}
rows=[]
start=time.monotonic()
for n in range(N):
    if n<4: an=win[n]
    else:
        P=ps(n-4)
        an,rem=divmod(-sum(P[i]*win[i] for i in range(4)),P[4])
        assert rem==0
        win=win[1:]+[an]
    if n<=100: assert an==direct(n)
    U+=an;T+=n*an;M2+=n*n*an;S+=(20*n+17)*an
    assert S%(n+1)==0,("counterexample",n+1)
    if n+1 in checkpoints:
        p,e=checkpoints[n+1]
        row={"N":n+1,"p":p,"e":e,"vU":val(U,p),"vT":val(T,p),"vS":val(S,p),"vM2":val(M2,p)}
        assert row["vU"]>=e and row["vS"]>=e
        if p!=5: assert row["vT"]>=e and row["vM2"]>=e-1
        rows.append(row)
    if (n+1)%10000==0:print("checked",n+1,"elapsed",round(time.monotonic()-start,2),flush=True)
result={"verified_through":N,"all_divisibilities_passed":True,"recurrence_divisions_exact":True,"direct_comparison_through":100,"checkpoints":rows,"seconds":round(time.monotonic()-start,3)}
print(json.dumps(result,indent=2),flush=True)
Path(__file__).with_suffix(".json").write_text(json.dumps(result,indent=2)+"\n")


from fractions import Fraction
R = Fraction(1)
for p in range(3,224,2):
    if not all(p%d for d in range(2,isqrt(p)+1)): continue
    e=0; power=1
    while power*p<=N: power*=p; e+=1
    good=(p==5 or p%20 in (11,13,17,19))
    nu=Fraction(1) if good else 1-Fraction(1,p**(e+1))
    R*=nu/(1-Fraction(1,p*p))
assert 8*R/Fraction(355,113)**2 > Fraction(9992,10000)
print("Exact density lower bound > 0.9992: PASS")

\end{lstlisting}
\begin{thebibliography}{99}
\raggedright
\bibitem{Sun2010}
Z.-W. Sun, \emph{On Ap\'ery numbers and generalized central trinomial
coefficients}, arXiv:1006.2776v11, 19 August 2010.
Definition of \(W_n(b,c)\), p.~29; Conjecture 5.8(i)--(ii), pp.~32--33.
\url{https://arxiv.org/abs/1006.2776v11}.

\bibitem{OEIS}
Z.-W. Sun, \emph{A179524}, The On-Line Encyclopedia of Integer Sequences,
entry dated 17 July 2010, conjecture attributed to 1 July 2010;
accessed 7 November 2026.
\url{https://oeis.org/A179524}.

\bibitem{SunRefinement}
Z.-W. Sun, \emph{A refinement of a congruence result by van Hamme and
Mortenson}, Illinois Journal of Mathematics \textbf{56} (2012), no.~3,
967--979. Theorem 1.2.
\url{https://arxiv.org/abs/1011.1902v6}.

\bibitem{Duke}
W. Duke, \emph{Modular functions and the uniform distribution of CM
points}, Mathematische Annalen \textbf{334} (2006), 241--252.
The first page gives \(j(i\sqrt5)=632000+282880\sqrt5\).
\url{https://math.ucla.edu/~wdduke/preprints/modud.pdf}.

\bibitem{EOY}
N. Elkies, K. Ono and T. Yang, \emph{Reduction of CM elliptic curves and
modular function congruences}, International Mathematics Research
Notices (2005), no.~44, 2695--2707.
The introduction recalls Deuring's nonsplit-prime reduction theorem.
\url{https://people.math.wisc.edu/~tonghaiyang/deuring.pdf}.

\bibitem{Silverman}
J. H. Silverman, \emph{The Arithmetic of Elliptic Curves}, second edition,
Graduate Texts in Mathematics 106, Springer, 2009.
Theorem V.4.1(a), p.~148.

\bibitem{Formal}
Google DeepMind and contributors, \emph{formal-conjectures},
\texttt{FormalConjectures/OEIS/179524.lean}, inspected 7 November 2026.
Initial path commit \texttt{1bb238f2be65a7d796fb29a15f1c89f764062515},
20 August 2026.
\url{https://github.com/google-deepmind/formal-conjectures/blob/main/FormalConjectures/OEIS/179524.lean}.

\bibitem{OEISOpen}
T. Adamczewski, \emph{OEIS Open: How many conjectures can language
models turn into theorems?}, arXiv:2608.11941v2, 13 August 2026.
\url{https://arxiv.org/abs/2608.11941v2}.
Released exact target:
\url{https://github.com/epoch-research/LeanOpenProblems/blob/main/apn/data/oeis/Isolated/oeis_a179524_sun_conjecture_2.lean}.

\bibitem{SunApery}
Z.-W. Sun, \emph{On sums of Ap\'ery polynomials and related congruences},
Journal of Number Theory \textbf{132} (2012), 2673--2699.
Theorem 1.1 and Remark 1.1.
\url{https://arxiv.org/abs/1101.1946v4}.

\bibitem{GuoZeng}
V. J. W. Guo and J. Zeng, \emph{Proof of some conjectures of Z.-W. Sun
on congruences for Ap\'ery polynomials}, Journal of Number Theory
\textbf{132} (2012), 1731--1740; arXiv:1101.0983.
Theorem 1.3.
\url{https://arxiv.org/abs/1101.0983}.

\bibitem{Babanskyy}
O. Babanskyy, \emph{Sun's supercongruence for two Ap\'ery-like binomial
sums at discriminant \(-20\)}, public manuscript, 2026.
Section 2, Lemmas 2.1--2.2, pp.~2--3.
Pinned repository commit
\texttt{4caad01fa469c86a1a95b930293354d843b79d64};
the inspected repository history dates from 3 October 2026.
\url{https://github.com/aconsciousfractal/Sun-3-7-Supercongruence/blob/4caad01fa469c86a1a95b930293354d843b79d64/paper/Sun_3_7_Supercongruence.pdf}.


\bibitem{Supplied}
QEEDDD Verifiable-Mathematical-Reasoning-Corpus project,
\emph{Sun's divisibility conjecture for
\(\sum_{k<n}(20k+17)W_k(1,-4)\): squarefree moduli, the second layer,
and a reduction via negative-index continuation},
supplied manuscript dated 7 November 2026, 13 pages.
The title page credits drafting with Claude (Anthropic);
the PDF and its single \TeX{} source were supplied for this study.

\bibitem{Delaygue}
\'E. Delaygue, \emph{Arithmetic properties of Ap\'ery-like numbers},
Compositio Mathematica \textbf{154} (2018), 249--274.
\href{https://doi.org/10.1112/S0010437X17007552}{doi:10.1112/S0010437X17007552}.
\url{https://arxiv.org/abs/1310.4131}.

\end{thebibliography}
\end{document}
